\documentclass[11pt]{article}
\usepackage[margin=0.9in]{geometry}
\usepackage{amsmath,amssymb,amsthm,booktabs,mathrsfs}
\usepackage[T1]{fontenc}
\usepackage{lmodern,microtype}
\usepackage[colorlinks=true,linkcolor=blue,urlcolor=blue]{hyperref}
\newtheorem{theorem}{Theorem}[section]
\newtheorem{lemma}[theorem]{Lemma}
\newtheorem{proposition}[theorem]{Proposition}
\newtheorem{corollary}[theorem]{Corollary}
\theoremstyle{remark}
\newtheorem{remark}[theorem]{Remark}
\newcommand{\norm}[1]{\left\|#1\right\|}
\title{Toward a four-dimensional Yang--Mills construction\\
\large Notebook 31: finite local induction and coupling extraction}
\author{Research working document}
\date{17 October 2026}
\begin{document}
\maketitle
\begin{abstract}
We carry the renormalized homogeneous local induction to actual
finite periodic regulators. A short-lift Maxwell extraction is
normalized by the actual finite classical action. The finite Ward
defects are repaired degree by degree; periodic summation by parts
and connected localization give a classical contraction with an
explicit exponentially small wrapping error. The carrier exponent
and field radius are unchanged.

The finite quantum map descends directly to equality classes of
normalized total interactions. Its input Lipschitz bound, its
quadratic zero-input force, and the finite classical contraction
give a uniform block induction. At every intermediate stage an
exactly equivalent connected family lies within the next integral's
strict budget. The change of representative introduces no
approximation error; all raw constants are compensated in the
vacuum factor. The number of permitted finite blocks does not
enter the amplitude constants.

We also prove evenness in the coupling at fixed retained threshold.
A uniform complex coupling disk and Cauchy estimates identify the
leading zero-input correction as the pivot term plus the Gaussian
determinant, with error bounded by
\(C[\tau(T)+g^2T^6]\). The physical logarithmic threshold is substituted
only after this two-parameter estimate is established.

On the jointly compatible infinite interaction space, a
quadratic-only spatial estimate at fixed field radius and the
classical semigroup prove exact preservation of the Maxwell
coefficient. Thus the coupling drift on an extracted input
comes entirely from the quantum remainder. This infinite
identity does not set finite cutoff-extraction errors to zero.

These results concern the specified local macrostep and its common
analytic continuation. The physical inverse coupling must satisfy
the exact extraction recurrence and remain in the weak range.
Its drift sign and ultraviolet tuning, the other field sectors,
and the continuum observable construction remain to be proved.
\end{abstract}

The goal remains a nontrivial four-dimensional Yang--Mills quantum
theory. This notebook uses the actual finite classical and quantum
constructions of Notebooks 15--30, with their fixed gauge, retained
window, trace convention, and common domains. It closes the finite
representative-budget issue left explicit in Notebook 30.



% No preamble. Prefix fcc-. Actual finite references and finite periodic backgrounds.
\section{Finite-regulator classical contraction}
\label{fcc-section}

Fix dimension four, a compact simple gauge group with the fixed
color normalization, a primitive odd ratio \(\ell\), and
\(0<\alpha<1\). A semisimple product uses one reference coefficient
per simple factor. Fix the activity exponent \(a_0>0\) and a
grouping \(J\). Let \(L\) be an allowed endpoint ratio and write
\begin{equation}
 P=LN.
 \label{fcc-periods}
\end{equation}
Here \(P\) is the input current-unit torus period and \(N\) is
the output current-unit period. Both cubic tori are compatible
with the partitions; in particular \(J\mid N\). The preceding
homogeneous depth \(M\) is arbitrary.

On a current torus of period \(Q\), let
\begin{equation}
 \|E\|_{Q;a_0,b}
   =\sup_i\sum_{H\ni i}e^{a_0|H|}\|E_H\|_b,
 \qquad
 \Phi_E(A)=\sum_H\bigl[E_H(A|_H)-E_H(0)\bigr],
 \label{fcc-norm-total}
\end{equation}
where \(H\) ranges over nonempty connected sets of \(J\)-cells
and \(\|\cdot\|_b\) is the supremum norm on the complex
scalar-coordinate \(b\)-polydisc. Quotient by equality of
\(\Phi_E\) on that full polydisc, and denote the quotient norm by
\(\|[E]\|_Q\). The equality kernel is closed, since finite-volume
evaluations are continuous, so this is a Banach quotient.
We use its invariant subspace: local gauge, translation,
hypercubic and global color identities are imposed on totals.
Individual coefficients need not be invariant. Vacuum constants
are tracked separately.

\subsection{Actual finite inputs and the reference projection}
\label{fcc-inputs}

The earlier finite constructions supply the following inputs.
First, for \(n=2,3,4\), fixed-order Cauchy estimates and connected
carriers give, for some fixed \(\mu=\mu(J)>0\),
\begin{equation}
 \sup_{b_1}\sum_{b_2,\ldots,b_n}
 e^{\mu D_Q(b_1,\ldots,b_n)}
 |K_{n,Q}(b_1,\ldots,b_n)|
 \le C_n(J)b^{-n}\|E\|_{Q;a_0,b},
 \qquad K_{n,Q}=D^n\Phi_E(0).
 \label{fcc-kernel-bound}
\end{equation}
Here \(D_Q\) is the maximum pairwise torus \(\ell^1\) distance
between bond midpoints, and internal indices are included in
the sums. Indeed, every differentiated tuple lies in a common
carrier, \(D_Q\le C J|H|\), while its fixed-degree coordinate
count is a polynomial in \(D(J)|H|\). Choosing \(\mu\) small
and spending the remaining cardinal reserve absorbs this
polynomial without a volume factor.

Second, \(S_{M,Q}\) is the actual finite homogeneous classical
reference. It has zero value and first derivative at zero and
the same total symmetries. Its strict reference family on one
fixed radius \(\rho\) is uniformly bounded by \(M_S\). Its flat
Hessian is the exact relative periodization of the uniformly
exponentially decaying \(q_M\), whose midpoint Maxwell
coefficient is one.

Removing coefficientwise homogeneous degrees zero and one
does not change its normalized total. The phase projection
bound gives, for \(0<b\le\rho/2\),
\begin{equation}
 \|S_{M,Q}^{(\ge2)}\|_{Q;a_0,b}
 \le \frac{M_S(b/\rho)^2}{1-b/\rho}
 \le C_S b^2,\qquad C_S=2M_S/\rho^2.
 \label{fcc-reference-small}
\end{equation}
This does not assert that the original individual reference
coefficients have zero linear jets.

For \(Q\ge64\), symmetrically lift the quadratic torus kernel
using midpoint diameter at most \(\lfloor Q/16\rfloor\), and
write its midpoint symbol as \(\Gamma^E_{2,Q}\). For distinct
directions \(\mu,\nu\), with the invariant color normalization,
put
\begin{equation}
 b_Q(E)=\frac12
       \partial_{p_\nu}^2\Gamma^E_{2,Q;\mu\mu}(0).
 \label{fcc-cutoff-jet}
\end{equation}
Point symmetry makes the distinct pair immaterial. The exact
flat reference periodization implies
\begin{equation}
 |b_Q(S_{M,Q})-1|\le C(J)e^{-c(J)Q}.
 \label{fcc-reference-denominator}
\end{equation}
This estimate includes both the discarded infinite kernel
tail and its long images folded into the short torus lift.
Uniformity of the \(q_M\) decay makes it uniform in \(M\).
Choose \(Q_0(J)\) so that the denominator has modulus at least
\(1/2\), and define
\begin{equation}
 \beta_{M,Q}([E])
     =\frac{b_Q(E)}{b_Q(S_{M,Q})},
 \qquad
 \mathcal P_{M,Q}[E]
     =[E]-\beta_{M,Q}([E])[S_{M,Q}].
 \label{fcc-projection}
\end{equation}
The kernels depend only on the normalized total, so these
are well defined on the quotient. Equations
\eqref{fcc-kernel-bound}--\eqref{fcc-reference-small} imply
\begin{align}
 |\beta_{M,Q}([E])|
     &\le C_\beta(J)b^{-2}\|[E]\|_Q,\notag\\
 \beta_{M,Q}([S_{M,Q}])&=1,\qquad
 \mathcal P_{M,Q}^2=\mathcal P_{M,Q},\qquad
 \|\mathcal P_{M,Q}\|\le C_{\rm ext}(J)
       :=1+C_\beta(J)C_S.
 \label{fcc-projection-bound}
\end{align}
The projection bound, unlike the bound on \(\beta\) alone,
is independent of the small \(b\).

Third, the actual finite pre-gauge ordinary representative
\(a=a_{M,L,N}(w)\) is the partial ordinary average of the
full \(ML\)-refined ordinary minimizer. The exact finite
two-minimum comparison gives
\begin{equation}
 a(0)=0,\qquad
 S_{M,LN}(a(w))=S_{ML,N}(w).
 \label{fcc-reference-transport}
\end{equation}
The identity first holds for real small fields and then
holomorphically. In dimension four it has no additional
dimensional prefactor.

For any fixed finite cardinal reserve \(A\), this map has
coherent complete-field support coefficients on final unit
cells:
\begin{align}
 a|_i&=\sum_{\varnothing\ne X}a_{i,X}(w_X),
       \qquad a_{i,X}(0)=0,\notag\\
 \sup_i\sum_X e^{\kappa\tau_N(\{i\}\cup X)+A|X|}
       \|a_{i,X}\|_{C^{1,\alpha};b}
      &\le C_* b/L.
 \label{fcc-pregauge-family}
\end{align}
The whole current-field output fiber and its physical
gradient H\"older quotients are inside the support sum.
Coefficients on overlapping fixed collars are restrictions
of one global assignment. The constants and the permitted
classical radius are uniform in \(M,L,N\). This is the
finite complete-field prefix theorem, not a periodization
assumption about an infinite background.

Define scalar classical substitution by
\begin{equation}
 (\mathcal T_{M,L,N}\Phi_E)(w)=\Phi_E(a_{M,L,N}(w)).
 \label{fcc-transport}
\end{equation}
For \(L\) sufficiently large, the amplitude is strictly
inside the input \(b\)-domain, so equality of finite totals
implies equality after substitution. The estimates below
give a bounded output representative.

We also record the finite-domain invariance argument.
For a real output field and a sufficiently small local
gauge transform, the finite transfer and gauge-flow
construction gives an actual near-identity connecting
input gauge \(g\). Both endpoint angles have size
\(\rho_0=C_A b/L\). The bond equation \(V^g=U\) bounds
the edge increments of \(g\) by
\(C(|U-I|+|V-I|)\), and hence
\(|D\log g|\le C\rho_0\) on its fixed near-identity chart.
The existence of the small chart comes from that actual
finite construction, not merely from an edge bound.
Every field along \(\exp(t\log g)\), \(0\le t\le1\), has
scalar-coordinate size at most \(C_g\rho_0\). Fixed
matrix/color norm conversions are included in \(C_g\).
Taking \(L>C_gC_A\) with a strict margin keeps the entire
path inside \(b\). The input Ward identity integrates
along it, and analytic continuation gives the local
complex output Ward identity. No infinite gauge
truncation is needed. The finite minimum comparisons
or exact covariance likewise give the spatial and
global color identities on this common domain.

Consequently \(\mathcal T\) acts on the finite invariant
quotients. The full-reference identity gives exactly
\begin{equation}
 \mathcal T_{M,L,N}[S_{M,LN}]=[S_{ML,N}],\qquad
 \mathcal P_{ML,N}\mathcal T_{M,L,N}\mathcal P_{M,LN}
       =\mathcal P_{ML,N}\mathcal T_{M,L,N}.
 \label{fcc-telescope}
\end{equation}
This does not assert that \(\mathcal T\) preserves the
zero-\(\beta\) subspace.

\subsection{Graded Ward defects and finite-range jet repair}
\label{fcc-repair-section}

Let \([F]\) satisfy \(\beta_{M,P}([F])=0\), and choose a
representative of norm at most \(\varepsilon\), as close
to its quotient norm as desired. Set \(T_0=\lfloor P/16\rfloor\).
The symmetric infinite lift
\begin{equation}
 \widetilde K_n(b_1,\ldots,b_n)
  =K_{n,P}(b_1\bmod P,\ldots,b_n\bmod P)
           {\bf1}_{\{D_\infty(b_1,\ldots,b_n)\le T_0\}}
 \label{fcc-short-lift}
\end{equation}
preserves slot symmetry, midpoint parity and the spatial
and color symmetries. A short torus tuple has a unique
lift with its first bond fixed: all coordinate differences
are below half a period, and the lifted pairwise distances
are the torus distances. Thus
\begin{equation}
 K_{n,P}=\operatorname{Per}_P\widetilde K_n+K_n^{\rm tail},
 \qquad
 K_n^{\rm tail}=K_{n,P}{\bf1}_{\{D_P>T_0\}}.
 \label{fcc-short-tail}
\end{equation}
Every fixed spatial moment of the tail is bounded by
\(C(J,m)\varepsilon b^{-n}e^{-c(J)P}\).

Ward operations are finite sums of nearest-neighbor
shifts, coincident-slot insertions and color commutators.
They commute with relative periodization. A lifted Ward
defect has diameter at most \(T_0+2\), where periodization
is injective. Its infinite fixed moments are therefore
bounded by those of the discarded tails.

Keep the field-degree grading. For \(b\le1\), the
degree-\(n\) identity involves only degrees \(2,\ldots,n\),
so its fixed-moment error is bounded by
\begin{equation}
 \delta_n=C(J)\varepsilon b^{-n}e^{-c(J)P},
       \qquad n=2,3,4.
 \label{fcc-graded-defect}
\end{equation}
The quadratic coefficient elimination uses only
\(\delta_2\); the cubic uses \(\delta_3\) and the already
bounded lower jets, and the quartic uses \(\delta_4\)
and its lower jets. Since \(b\le1\), these lower terms
fit the displayed grading. A common \(b^{-4}\) estimate
for all degrees is neither needed nor used.

For completeness, the lifted hypercubic Hessian has
mass \(m\) and quadratic coefficients \(a,b,c\).
The degree-one and degree-three Ward comparisons bound
\(m,a,c\) by \(C\delta_2\), because the extracted
coefficient \(b_P(F)\) is exactly zero. Odd midpoint
orders vanish by parity. Write the cubic linear polynomial
as \(A p+B q\); the independent \(p_\mu q_\sigma\) Ward
coefficients bound \(B\), and the \(q_\nu p_\sigma\)
coefficients bound \(A\), by \(C\delta_3\). The quartic
leading identity contracts its constant tensor with an
arbitrary gauge momentum and bounds it by
\(C\delta_4\). Hence the lifted quadratic jets through
degree three, cubic jets through degree one, and quartic
constant are respectively \(O(\delta_2),O(\delta_3),
O(\delta_4)\).

Pass to tail-coordinate Laurent symbols \(F_n(z)\),
anchoring the first positive-bond tail and retaining
internal indices. Fixed midpoint phases and the local
change \(z=e^{ip}\) preserve these finite-order bounds.
Put \(m_2=3,m_3=1,m_4=0\), and take
\begin{equation}
 J_{m_n}F_n(z)=
    \sum_{|\gamma|\le m_n}
       \frac{\partial_z^\gamma F_n(1)}{\gamma!}
                         (z-1)^\gamma.
 \label{fcc-laurent-jet}
\end{equation}
Its kernel is of fixed finite relative range and has
absolute row norm at most \(C\delta_n\).

Average this kernel over the \(n!\) slot permutations,
reanchoring the first slot each time, and call the result
\(R_n\). Its relative range increases only by a fixed
factor and its norm remains \(O(\delta_n)\).
The corrected kernel
\begin{equation}
 \widehat K_n=\widetilde K_n-R_n
 \label{fcc-repaired-kernel}
\end{equation}
has exactly zero jets through degree \(m_n\). Indeed
each reanchored residual is the composition of
\(F_n-J_{m_n}F_n\) with an analytic monomial change
fixing \(z=1\), possibly times a harmless analytic
phase. This preserves vanishing order \(m_n+1\);
averaging preserves it as well. Reality and color
covariance can be retained. Point symmetry of the
repair is not required by the spatial estimate.

For a fixed smaller \(\sigma=\sigma(J)>0\),
\begin{align}
 \sup_{\text{first index}}\sum_{\text{others},\,r_2,\ldots,r_n}
 e^{\sigma\sum_{j\ge2}|r_j|}
      |\widehat K_n(0,r_2,\ldots,r_n)|
       &\le C(J)\varepsilon b^{-n},\notag\\
 \|R_n\|_{\rm row}
       &\le C(J)\varepsilon b^{-n}e^{-c(J)P}.
 \label{fcc-repaired-bounds}
\end{align}
The second line has fixed finite relative range.
The finite polynomial decomposition is exact:
\begin{equation}
 \frac{K_{n,P}(a,\ldots,a)}{n!}
  =\frac{(\operatorname{Per}_P\widehat K_n)(a,\ldots,a)}{n!}
   +\frac{(\operatorname{Per}_P R_n)(a,\ldots,a)}{n!}
   +\frac{K_n^{\rm tail}(a,\ldots,a)}{n!}.
 \label{fcc-polynomial-decomposition}
\end{equation}
No error term is discarded or separately asserted to
be gauge invariant.

\subsection{Direct periodic spatial localization}
\label{fcc-periodic-spatial-section}

We give the finite version of the spatial argument explicitly.
For a relative infinite kernel \(\widehat K_n\), its
periodized polynomial equals
\begin{equation}
 \frac1{n!}\sum_{x\in T_P}
  \sum_{r_2,\ldots,r_n\in\mathbb Z^4}
   \widehat K_n(r_2,\ldots,r_n)
       [a(x),a(x+r_2),\ldots,a(x+r_n)],
 \label{fcc-periodic-polynomial}
\end{equation}
where fields are read periodically. The relative sums
converge absolutely, and repeated torus slots retain
their multilinear multiplicity.

For the quadratic symbol \(F(z)=\sum_r\widehat K_2(r)z^r\),
define
\[
 B_jF(z)=
 \frac{F(1,\ldots,1,z_j,\ldots,z_d)
       -F(1,\ldots,1,1,z_{j+1},\ldots,z_d)}{z_j-1}.
\]
The quotient is removable, and
\(F-F(1)=\sum_j(z_j-1)B_jF\). A monomial contributes
\(|r_j|\) terms of exponent length at most \(|r|_1\).
Thus the combined row Wiener norm, including the \(j\)
sum, loses at most \([e(\sigma-\sigma')]^{-1}\).
Vanishing order decreases by at most one. Apply twice
and put \(\widetilde H_{\rho\tau}=-z_\rho B_\tau B_\rho F\).
Then
\begin{equation}
 F=\sum_{\rho,\tau}(z_\rho^{-1}-1)
          \widetilde H_{\rho\tau}(z)(z_\tau-1),
 \qquad \widetilde H_{\rho\tau}(1)=0.
 \label{fcc-divergence-factor}
\end{equation}
The combined row norm at exponent \(\sigma/2\) is at
most
\[
 \frac{16e^{\sigma/2}}{e^2\sigma^2}
                  C(J)\varepsilon b^{-2}.
\]

Finite periodic summation by parts has no boundary term.
Equation \eqref{fcc-divergence-factor} and the zero sum
of \(\widetilde H\) give exactly
\begin{equation}
 \frac12\sum_{\substack{x\in T_P,\,r\\\rho,\tau}}
 [D_\rho a(x)]^T\widetilde H_{\rho\tau}(r)
                [D_\tau a(x+r)-D_\tau a(x)].
 \label{fcc-periodic-ibp}
\end{equation}
No summability of an infinite field is assumed.

For the cubic term set
\(\ell_r(x)=\sum_\rho r_\rho D_\rho a(x)\) and
\(\operatorname{Rem}_r=a(x+r)-a(x)-\ell_r(x)\).
The exact lifted constant and first moments cancel
the all-anchor and single-\(\ell\) terms. The six
remaining terms are the two \(a,a,\operatorname{Rem}\)
terms, \(a,\ell,\ell\), the two
\(a,\ell,\operatorname{Rem}\) terms, and
\(a,\operatorname{Rem},\operatorname{Rem}\).
For the quartic term subtract the all-anchor product
and telescope its three shifted factors. These are
identities for periodic fields with lifted integer
offsets; no periodic affine function is introduced.

Let \(h=L^{-1}\) and \(m=\sum_{j\ge2}|r_j|_1\).
For \(m\le L\), coherence in
\eqref{fcc-pregauge-family} gives, in the weighted
family norm on one fixed enlarged collar,
\begin{align}
 |a|&\le Cbh,& |Da|&\le Cbh^2,\notag\\
 |Da(x+r)-Da(x)|&\le Cbh^{2+\alpha}|r|^\alpha,&
 |\operatorname{Rem}_r|&\le
                 Cbh^{2+\alpha}|r|^{1+\alpha}.
 \label{fcc-near-estimates}
\end{align}
Negative path traversals use the positive difference
at the lower endpoint. Products of coefficient families
glue their support trees and preserve cardinal reserve.
The quadratic, cubic and quartic residual products
are therefore bounded respectively by
\[
 Cb^2h^{4+\alpha}|r|^\alpha,\qquad
 Cb^3h^{4+\alpha}(1+m)^6,\qquad
 Cb^4h^5(1+m).
\]
The kernel moments cost only fixed powers of
\(\sigma^{-1}\), and their \(b^{-n}\) cancels \(b^n\).

For every product, project the finite endpoint root
set and every terminal set \(\{\text{root}\}\cup X\)
to final \(J\)-cells. Choose the short connecting
graphs supplied by the projection lemma and take
their connected union \(H\). The carrier need not
contain the full projection of an arbitrarily chosen
fine tree. With a fixed number \(n_0\) of endpoint
roots,
\begin{align}
 |H|&\le1+T_H,\notag\\
 T_H&\le\frac2J\left[
          \tau_N(\text{roots})
           +\sum_{\rm occ}\tau_N(\{\text{root}\}\cup X)\right]
             +4\left[n_0+\sum_{\rm occ}|X|\right].
 \label{fcc-carrier-geometry}
\end{align}
Here \(\tau_N\) uses final unit cells of period \(N\).
Shortest periodic paths obey the same bound; wrapping
can only reduce the union cardinality. If \(j(x)\)
is the final \(J\)-cell containing \(x/L\), then
\begin{equation}
 e^{a_0|H|}{\bf1}_{\{z\in H\}}
 \le e^{a_0}e^{-\operatorname{dist}_N(z,j(x))}
                         e^{(a_0+1)T_H}.
 \label{fcc-reroot}
\end{equation}
The distance here is on final \(J\)-cells, of period
\(N/J\), not on unit cells. Choose
\begin{equation}
 A\ge4(a_0+1),\qquad 2(a_0+1)/J\le\kappa.
 \label{fcc-reserve}
\end{equation}
These pay the support and attached-tree weights.
Finally,
\begin{equation}
 \sum_{x\in T_P}e^{-\operatorname{dist}_N(z,j(x))}
       \le(JL)^4 S_1,\qquad
 S_1=\sum_{y\in\mathbb Z^4}e^{-|y|_1}.
 \label{fcc-anchor-count}
\end{equation}
There are exactly \((JL)^4\) current anchor sites per
output \(J\)-cell. Choosing a shortest representative
per torus coset bounds the remaining sum by \(S_1\).

For \(m>L\), use separate endpoint roots. Their
joining tree has length at most \(C(m/L+1)\), so
\eqref{fcc-reroot} adds at most
\(C e^{C'm/(JL)}\). The rough quadratic, cubic and
quartic residual bounds are respectively
\(Cb^2h^4\), \(Cb^3h^3(1+m)^2\), and \(Cb^4h^4\).
Taking \(L\) large after \(J\) makes \(C'/(JL)\le\sigma/8\).
The unspent kernel decay bounds all far terms by
\[
 C(J)\varepsilon(L+2)e^{-\sigma L/8}
       \le C'(J)\varepsilon L^{-\alpha}.
\]
The near terms have the same or better bound after
\eqref{fcc-anchor-count}. Normal regrouping by \(H\)
therefore constructs a connected output family with
\begin{equation}
 \left\|\sum_{n=2}^4
       \frac{(\operatorname{Per}_P\widehat K_n)
                       (a,\ldots,a)}{n!}\right\|_{\rm rep}
       \le C_{\rm low}(J)\varepsilon L^{-\alpha}.
 \label{fcc-canceled-estimate}
\end{equation}
The notation \(\|\cdot\|_{\rm rep}\) denotes the norm
of the constructed representative, hence bounds its
quotient norm. Its coefficients are holomorphic on
the same \(b\)-polydisc. This proof used the actual
finite periodic support family, not a putative
infinite-volume extension of it.

\subsection{Wrapping errors at the unchanged activity exponent}
\label{fcc-errors-section}

We use a direct endpoint-product bound. Suppose a torus
\(n\)-kernel \(W\), \(n\le4\), satisfies
\begin{equation}
 \sup_{b_1}\sum_{b_2,\ldots,b_n}e^{\mu D_P}|W|
                         \le A_n.
 \label{fcc-error-kernel}
\end{equation}
Join each endpoint to the first by a shortest torus
path. Its root tree on final unit cells has length
at most \(C_n(D_P/L+1)\). Insert the actual endpoint
support families and use
\eqref{fcc-carrier-geometry}--\eqref{fcc-reroot}.
The additional factor is at most
\[
 C\exp\!\left(\frac{C_n'(a_0+1)D_P}{JL}\right).
\]
Increase the lower threshold for \(L\) until this
exponent coefficient is at most \(\mu/4\).
The field product contributes \((Cb/L)^n\), while
\eqref{fcc-anchor-count} gives
\begin{align}
 \left\|\frac{W(a,\ldots,a)}{n!}\right\|_{\rm rep}
   &\le C(J,n)A_n b^n L^{4-n},\notag\\
 \left\|\frac{(W{\bf1}_{\{D_P>T_0\}})(a,\ldots,a)}
                                      {n!}\right\|_{\rm rep}
   &\le C(J,n)A_n b^n L^{4-n}e^{-\mu T_0/2}.
 \label{fcc-tail-localization}
\end{align}
The target exponent is exactly \(a_0\). The extra
carrier length was paid after spatial compression
by \(JL\), rather than by lowering the old-scale
activity exponent.

Apply the second inequality to \(K_n^{\rm tail}\),
using \(A_n\le C(J)\varepsilon b^{-n}\) and
\(T_0\ge P/32\). Apply the first to the fixed-range
repair \(R_n\), retaining its factor \(e^{-c(J)P}\)
from \eqref{fcc-repaired-bounds}. Relative periodization
counts every occurrence, including any repeated
periodic slots. Together these give
\begin{equation}
 \left\|\sum_{n=2}^4
    \frac{(\operatorname{Per}_P R_n+K_n^{\rm tail})
                         (a,\ldots,a)}{n!}\right\|_{\rm rep}
 \le C_{\rm err}(J)\varepsilon
              (L^2+L+1)e^{-c(J)LN}.
 \label{fcc-error-estimate}
\end{equation}
All \(b\)-powers cancel in their correct degree.
An old-scale tail estimate that loses part of
\(a_0\) is not being invoked.

\subsection{High degrees and the finite quotient estimate}
\label{fcc-contraction-section}

Coefficientwise homogeneous phase projection has norm
at most one. Thus
\begin{equation}
 R_{\ge5}E=E-\sum_{n=0}^4P_nE,\qquad
 \|R_{\ge5}E\|\le6\|E\|.
 \label{fcc-high-projection}
\end{equation}
The supported fifth-order background theorem gives
\begin{equation}
 \|\mathcal T_{\rm high}R_{\ge5}E\|_{N;a_0,b}
        \le6C_{\rm high}\varepsilon/L,
 \qquad C_{\rm high}=C_0(8C_A)^5.
 \label{fcc-high-estimate}
\end{equation}
After the fixed classical reserves, \(C_{\rm high}\)
does not depend on \(M,L,J,N\). A fixed \(J\) satisfying
the earlier tail budgets can be used while increasing
the endpoint to \(L=R^q\): the core-compression
coefficient decreases, whereas the \(D(J)t\) and
projection budgets are unchanged. The field-scaling
disk has radius \(L/(8C_A)\); choose \(b\) afterward
so that
\begin{equation}
 b\le1,\quad b\le\rho/2,\quad
 bL/(8C_A)\le r_A,\quad L/(8C_A)\ge2,
 \label{fcc-radius}
\end{equation}
and impose the other fixed classical domain
restrictions.

The degree-zero sector disappears from the normalized
total. The degree-one total vanishes by semisimple
global color invariance and remains zero after
substitution, even if its individual coefficients
are nonzero. The degree-two through degree-four
polynomials are exactly
\eqref{fcc-polynomial-decomposition}. Thus the low
estimate, error estimate, and high estimate represent
the entire finite normalized substituted total,
not merely its first few kernels.

\begin{theorem}[Finite classical contraction]
\label{fcc-theorem}
Fix \(J\) and its stated classical tail and support
budgets. There are \(L_0(J)\), \(N_0(J)\), positive
\(c(J)\), and constants \(C_{\rm low}(J)\),
\(C_{\rm err}(J)\), \(C_{\rm ext}(J)\), independent
of \(L,b,M,N\), such that every allowed
\(L\ge L_0(J)\), every compatible \(N\ge N_0(J)\),
and every sufficiently small \(b\) chosen after \(L\)
satisfy
\begin{align}
 &\|\mathcal P_{ML,N}\mathcal T_{M,L,N}
                              \mathcal P_{M,LN}\|\notag\\
 &\quad\le C_{\rm ext}(J)^2
       \left[
       C_{\rm low}(J)L^{-\alpha}
       +\frac{6C_{\rm high}}L
       +C_{\rm err}(J)(L^2+L+1)e^{-c(J)LN}\right].
 \label{fcc-main-bound}
\end{align}
The input and output finite invariant quotients have
the same exponent \(a_0\) and radius \(b\).
Separate input and output projection constants can
replace the square if their conventions differ.
\end{theorem}

\begin{proof}
Take an almost-minimal representative of the extracted
class. Its norm can approach
\(\|\mathcal P_{M,LN}[E]\|\le C_{\rm ext}(J)\|[E]\|\).
Apply \eqref{fcc-canceled-estimate},
\eqref{fcc-error-estimate}, and
\eqref{fcc-high-estimate} to this representative.
They give an exactly equivalent connected output
family. The output projection costs the second
\(C_{\rm ext}(J)\). Take the quotient infimum and
send the representative error to zero.
\end{proof}

Renaming constants, the bound has the form
\begin{equation}
 \|\mathcal P_{\rm out}\mathcal T\mathcal P_{\rm in}\|
 \le A(J)L^{-\alpha}+B(J)/L
                       +E(L,J)e^{-c(J)N}.
 \label{fcc-renamed-bound}
\end{equation}
The sharper \eqref{fcc-main-bound} displays the actual
input period \(LN\). For any desired \(\theta>0\),
choose a finite allowed power \(L=R^q\) after the
fixed \(J\) to make the first two terms at most
\(\theta/2\); enlarge the compatible threshold
\(N_0\) until the error is at most \(\theta/2\) for
all \(N\ge N_0\); and choose \(b\) afterward.
This is a \(\theta\)-contraction uniform in preceding
depth and all these finite regulators. Shrinking
\(b\) does not change its constants.

By \eqref{fcc-telescope} the same bound applies to
\(\mathcal P_{\rm out}\mathcal T\), with the input
extraction implicit. That identity uses the full
actual finite reference, not merely an agreement
of flat Hessians.

\subsection{Scope}
\label{fcc-scope}

The estimate uses actual finite normalized totals
and periodic spatial algebra. It requires neither
an infinite-volume interaction limit nor evaluation
of an infinite total at a nonsummable periodic
field. Small finite lifted Ward defects and
wrapping tails are retained explicitly. The
resulting small representatives are exactly
equivalent on the full finite \(b\)-polydisc; the
original carrier assignment need not itself
obey the improved norm.

The order is fixed \(J\), then a finite endpoint
\(L=R^q\), then the small \(b\). This removes the
finite-regulator classical relocalization issue.
It does not identify \(q\) sequential quantum
cutoff integrals with one direct macro cutoff
integral. Application to an actual finite quantum
recursion still needs the exact finite scalar
descent, every intermediate analytic and Gaussian
budget, the running-coupling bounds and the
specified cutoff/history identities. Raw vacuum
normalizations remain separate. No beta-function
sign, indefinitely iterable estimate, large-field
control or continuum existence claim follows
from this classical result alone.


% No preamble. Prefix fqi-. Cross-references only to this fragment and fcc-.
\section{Exact finite local quantum induction}
\label{fqi-section}

This section combines the actual common-domain finite macro integral,
finite-total quotient descent, the fixed-outer-radius force of
Notebook30, and Theorem~\ref{fcc-theorem}. The result concerns the
selected homogeneous \emph{local} macrostep and its specified
holomorphic continuation. It does not identify the full Wilson
pushforward, other field sectors, or arbitrary mixed histories.

\subsection{Finite geometry, totals and the actual map}
\label{fqi-setup-section}

Fix the primitive ratio $\ell$, an elementary composite ratio
$R=\ell^s$, auxiliary grouping $J$, and exponent $a_0=1/16$
satisfying the strict common-domain budgets. At stage $j$ let
\begin{equation}
 M_j=M_0R^j,\qquad Q_j=Q_0/R^j.
 \label{fqi-scales}
\end{equation}
The preceding homogeneous classical depth is $M_j$; the cubic
current-unit torus period is $Q_j$, after resetting spacing to one.
All periods are compatible integers and $J\mid Q_j$. Constants
may depend on the fixed geometry and color conventions, but are
uniform in depth and in the permitted periods above a lower
threshold.

On the period-$Q$ torus use the finite activity space and quotient
of \eqref{fcc-norm-total}. For clarity, the raw and normalized totals
are
\begin{equation}
 E^{\rm raw}(A)=\sum_HE_H(A|_H),\qquad
 \Phi_E(A)=E^{\rm raw}(A)-E^{\rm raw}(0),\qquad
 \|E\|_{Q;a_0,b}=\sup_i\sum_{H\ni i}e^{a_0|H|}\|E_H\|_b.
 \label{fqi-totals}
\end{equation}
There are finitely many coordinates and connected $J$-cell carriers
at each regulator. Equality of $\Phi_E$ on the full product ball
is a closed linear kernel, so the quotient is Banach. We impose
local gauge, translation, hypercubic and global color identities
on totals, not individual coefficients. Use a simple Lie algebra
for the scalar Maxwell functional; a finite componentwise version
handles a semisimple product.

Let $S_{M,Q}$ be the actual finite homogeneous ordinary minimum
action, with $S_{M,Q}(0)=0$. It is not replaced by a canonical
periodization of an independently selected infinite reference.
The actual common-domain construction supplies
\begin{align}
 \mathcal F_j(E)&=\mathcal T_jE+\mathcal G_j(E),\notag\\
 \mathcal G_j(E)&=\Delta_j+\mathcal Q_{{\rm det},j}
                                      +\mathcal L_j(Q_E).
 \label{fqi-map}
\end{align}
Here $\mathcal T_j$ is pre-gauge substitution, $Q_E$ is the exact
inherited fluctuation difference, and $\mathcal L_j$ is the
strict Gaussian logarithm with the fixed actual base exponent,
precision, retained window and chart. The coupling $g_j$ and
threshold $T_{{\rm cut},j}$ are held fixed when comparing inputs
to this one map. Field-independent factors are tracked separately.

On the original real common overlap, this is the exact chosen
local macro integral. On the common radius-$b$ ball it is its
proved holomorphic continuation, evaluated through pre-gauge
fields inside the old activity domain. Identification throughout
that ball with a globally defined gauge-orbit density requires
the extra orbit-representation hypothesis of the common-domain
theorem. It is not inferred from a quotient identity.

\subsection{Exact finite quotient descent}
\label{fqi-descent-section}

\begin{lemma}[Finite descent of the actual quantum update]
\label{fqi-descent-lemma}
The actual finite maps $\mathcal F_j$, $\mathcal T_j$ and
$\mathcal G_j$ descend to the invariant finite quotient on
their admissible ball. No infinite periodization, kernel limit
or translation averaging is required.
\end{lemma}
\begin{proof}
If $[E]=[E']$, then on the entire input product ball
\begin{equation}
 E^{\rm raw}(A)-E'^{\rm raw}(A)=c,
 \qquad c=E^{\rm raw}(0)-E'^{\rm raw}(0).
 \label{fqi-constant-difference}
\end{equation}
The actual pre-gauge field $a_p(w)$ and all fluctuation fields
$f_\lambda(w,\xi)$ used in the inherited construction lie in
that ball on the fixed retained window and the established
interpolation domains. Consequently
\begin{equation}
 E^{\rm raw}(f_1)-E^{\rm raw}(f_0)
       =E'^{\rm raw}(f_1)-E'^{\rm raw}(f_0).
 \label{fqi-same-insertion}
\end{equation}
The scalar inherited Gaussian insertion is identical, while
the exterior background terms differ by $c$. The covariance,
base action, chart, sharp window and measure factors are unchanged.
The exact finite integrals therefore differ only by $e^c$.
Their normalized holomorphic logarithms agree on the connected
domain, using the common nonvanishing branch. A raw logarithm
constant belongs to the vacuum normalization and disappears
from the quotient. The localized arrays need not agree.

Classical substitution descends immediately from
\eqref{fqi-constant-difference} evaluated at $a_p(w)$;
subtracting it gives descent of $\mathcal G_j$.
The exact real-sector symmetry and holomorphic Ward continuation
preserve the stipulated total identities. No carrierwise gauge
identity is assumed.
\end{proof}

The constant in \eqref{fqi-constant-difference} must be compensated
in a raw density, as specified below. For a real class, a chosen
representative can be made real without increasing its norm by
averaging $E_H(w)$ with $\overline{E_H(\overline w)}$.

\subsection{Uniform finite difference and forcing estimates}
\label{fqi-estimates-section}

Let $\bar\epsilon$ be the family ceiling of the strict two-input
comparison. For quotient inputs $x,y$ of norm at most
$\bar\epsilon/4$, choose near-minimal representatives $Y$ of $y$
and $D$ of $x-y$, and put $X=Y+D$. For arbitrarily small
$\delta>0$,
\begin{align}
 \|Y\|&\le\bar\epsilon/4+\delta,\notag\\
 \|D\|&\le\|x-y\|+\delta\le\bar\epsilon/2+\delta,\notag\\
 \|X\|&\le3\bar\epsilon/4+2\delta<\bar\epsilon.
 \label{fqi-compatible-representatives}
\end{align}
Both complete inputs are admissible and have invariant totals.
Covariance of their arrays is unnecessary at finite volume.

The actual Gaussian interpolation retains the base exponent and
precision. In particular impose
\begin{equation}
 v\le\frac{\eta_*}{64e},\qquad
 a_Q\bar\epsilon\le\frac{\eta_*}{128e},\qquad
 a_Q=16C_TC_R\frac{g_jT_{{\rm cut},j}}b.
 \label{fqi-comparison-budgets}
\end{equation}
The inherited difference is linear in $D$. The finite
Banach-valued logarithm Cauchy estimate gives
\begin{equation}
 \|\mathcal G_j(x)-\mathcal G_j(y)\|_{Q_{j+1}}
 \le L_{\mathrm{quant},j}\|x-y\|_{Q_j},\qquad
 L_{\mathrm{quant},j}
   =C_{\mathrm{quant}}\frac{g_jT_{{\rm cut},j}}b,\quad
 C_{\mathrm{quant}}=1024eS_1C_TC_R.
 \label{fqi-quantum-Lipschitz}
\end{equation}
Apply the family bound to the compatible $X,Y$, use
Lemma~\ref{fqi-descent-lemma}, and let $\delta$ decrease to zero.
No Fatou or local-limit argument is needed at a fixed regulator.
The external determinants cancel in this difference, and no
equal-low-jet hypothesis is used.

For the zero input, choose one outer field radius $\rho$ after
$R,J$ with all strict classical and Gaussian margins. Fix the
same actual threshold, charts, covariance and branch before
restricting to $b$. Notebook30's finite outer-domain force has
a uniform anchored bound $H_*$ on $\rho$. Its total linear
jet vanishes by semisimple adjoint invariance. Removing
coefficientwise degrees zero and one and using homogeneous
phase projection gives, for $b\le\rho/2$,
\begin{equation}
 \|\mathcal G_j(0)\|_{Q_{j+1}}\le C_Fb^2,
 \qquad C_F=2H_*/\rho^2.
 \label{fqi-force}
\end{equation}
This is an exactly equivalent normalized representative, not
a low-jet approximation. Its vacuum constant is kept separately.
The full force need not be $O(g^2)$: the pivot and Gaussian
determinant terms can persist at zero coupling.

\subsection{Finite references and the classical block}
\label{fqi-classical-section}

At stage $j$, use the finite Maxwell ratio of
\eqref{fcc-projection}, normalized against the actual reference
$S_j=[S_{M_j,Q_j}]$. For periods above its fixed threshold,
\begin{equation}
 \beta_j(S_j)=1,\qquad
 \|\beta_j\|\le C_\beta b^{-2},\qquad
 P_j=I-S_j\beta_j,\qquad
 \|P_j\|\le p:=1+C_\beta C_S,\qquad
 \|S_j\|\le C_Sb^2.
 \label{fqi-projections}
\end{equation}
The $b^{-2}$ factor is retained and cancels only in the
combined projection estimate. All constants are uniform in
the permitted periods and preceding depth.

The actual finite two-minimum comparison and ordinary
composition give
\begin{equation}
 \mathcal T_jS_j=S_{j+1},\qquad
 P_{j+1}\mathcal T_jP_j=P_{j+1}\mathcal T_j.
 \label{fqi-reference-telescope}
\end{equation}
In dimension four the fine action has the identical
$\eta^{d-4}=1$ normalization in both unit conventions.
The identity concerns the full actual action, not merely
its quadratic symbol. Put $A_j=P_{j+1}\mathcal T_j$.
The finite classical semigroup for invariant totals yields
\begin{equation}
 A_{j+q-1}\cdots A_jP_j
     =P_{j+q}\mathcal T_{M_j,R^q,Q_{j+q}}P_j.
 \label{fqi-block-telescope}
\end{equation}
The direct map on the right is classical substitution.
It does not replace the intervening quantum integrals or
their sharp windows by a one-shot integral.

Write $L=R^q$. Theorem~\ref{fcc-theorem} gives the block bound
\begin{equation}
 \theta(q,N):=p^2\left[
 C_{\mathrm{low}}(J)L^{-\alpha}+\frac{6C_{\mathrm{high}}}L
 +C_{\mathrm{err}}(J)(L^2+L+1)e^{-c(J)LN}\right],
 \qquad N=Q_{j+q}.
 \label{fqi-classical-block-bound}
\end{equation}
It retains the graded finite Ward defects and wrapping tails
at the same $a_0,b$. Fix $J$, choose a finite $q$, and then
a compatible $N_0$ above all reference, collar and period
thresholds so that
\begin{equation}
                 \theta(q,N)\le1/4\qquad(N\ge N_0).
 \label{fqi-period-threshold}
\end{equation}
Only afterward choose $b$ within the direct $L$ requirements
of \eqref{fcc-radius} and all common domains. One-step
transport has uniform bounds
\begin{equation}
 \|\mathcal T_j\|\le B_T,\qquad
 \|A_j\|\le B,\qquad B\ge\max(1,pB_T).
 \label{fqi-one-step-bound}
\end{equation}
The finite projection telescope remains exact although the
period and functional $\beta_j$ change at each stage.
It requires no identity $\beta_{j+1}\mathcal T_j=\beta_j$.

\subsection{A uniform finite block envelope}
\label{fqi-envelope-section}

Define the extracted maps
\begin{equation}
 \Psi_j(x)=P_{j+1}\mathcal F_j(x)=A_jx+H_j(x),\qquad
 H_j=P_{j+1}\mathcal G_j.
 \label{fqi-extracted-map}
\end{equation}
Let $\nu$ bound $pL_{\mathrm{quant},j}$ over the step data.
Equations \eqref{fqi-quantum-Lipschitz}--\eqref{fqi-force} give
\begin{equation}
 \operatorname{Lip}H_j\le\nu,\qquad
 \|H_j(0)\|\le C_0b^2,\qquad C_0=pC_F.
 \label{fqi-perturbation-bounds}
\end{equation}
Set
\begin{equation}
 K=B+1,\qquad L_q=\sum_{l=0}^{q-1}K^l,\qquad
 K_{\mathrm{amp}}=\max\{1,2C_0L_q\},\qquad
 C_{\mathrm{int}}=K^qK_{\mathrm{amp}}+C_0L_q.
 \label{fqi-envelope-constants}
\end{equation}
Choose the common $b$ so that $C_{\mathrm{int}}b^2<\bar\epsilon/4$,
and finally restrict the weak-coupling range so every strict
budget holds and
\begin{equation}
 \nu\le\min\{1,(4qK^{q-1})^{-1}\}.
 \label{fqi-nu-choice}
\end{equation}
For example, a sufficiently large fixed multiple of
$\sqrt{\log(1/g)}$ with a fixed lower threshold has
$gT(g)/b\to0$, $gT(g)^3\to0$ and $\tau(T(g))\to0$.
The outer force remains uniform on this smaller coupling range.

\begin{theorem}[Finite local induction with exact small representatives]
\label{fqi-theorem}
Use the actual finite maps, references and strict domains above.
Choose $q,N_0,b$ and the weak-coupling range in the stated order.
For every compatible finite history of $n$ steps with
\begin{equation}
 Q_0=R^nQ_n,\qquad Q_n\ge N_0,
 \label{fqi-history-periods}
\end{equation}
and initial extracted class satisfying
\begin{equation}
 \beta_0(x_0)=0,\qquad
 \|x_0\|\le r_b:=K_{\mathrm{amp}}b^2,
 \label{fqi-initial-ball}
\end{equation}
the following conclusions hold for couplings in the permitted
weak range.

Every full $q$-step block preserves the radius-$r_b$ extracted
ball and has Lipschitz constant at most $1/2$ when compared
with the same prescribed coupling data. Every intermediate
class has norm at most $C_{\mathrm{int}}b^2<\bar\epsilon/4$.
At every stage one can choose an exactly equivalent connected
finite representative on the same $a_0,b$ with norm at most
$2C_{\mathrm{int}}b^2<\bar\epsilon/2$. Thus every actual next
input and Gaussian budget remains valid. The changes of raw
constant are compensated by the exact vacuum factors below.

The constants are independent of the permitted number of
blocks and of preceding depth. The result describes the
selected homogeneous local germ and its continuation; its
global measure interpretation retains the stated extra
orbit-domain requirements.
\end{theorem}

\begin{proof}
For a block starting at stage $j$ with an input in
\eqref{fqi-initial-ball}, the one-step bound gives, inductively
for $0\le k\le q$,
\begin{equation}
 \|x_{j+k}\|
 \le K^kK_{\mathrm{amp}}b^2+C_0b^2\sum_{l<k}K^l
 \le C_{\mathrm{int}}b^2<\bar\epsilon/4.
 \label{fqi-intermediate-bound}
\end{equation}
This verifies the next map's domain before that map is invoked.
It also applies to the zero-started comparison trajectory.

For two trajectories using the same couplings, variation of
constants bounds their difference by
\begin{align}
 \|\delta_q\|
 &\le\left[\theta(q,Q_{j+q})
       +\nu\sum_{k=0}^{q-1}B^{q-1-k}(B+\nu)^k\right]
                                             \|\delta_0\|\notag\\
 &=\left[\theta(q,Q_{j+q})+(B+\nu)^q-B^q\right]
                                             \|\delta_0\|\notag\\
 &\le\left[\frac14+q\nu K^{q-1}\right]\|\delta_0\|
 \le\frac12\|\delta_0\|.
 \label{fqi-Duhamel}
\end{align}
Only the initial full classical product uses
\eqref{fqi-classical-block-bound}; every shorter product
uses powers of $B$. The block output at zero has norm at
most $C_0L_qb^2\le r_b/2$. Comparing with that output
therefore proves invariance of the radius-$r_b$ ball.
The output is extracted at every step by definition.

Repeat by full blocks; any final partial block satisfies
\eqref{fqi-intermediate-bound}. Condition
\eqref{fqi-history-periods} ensures every full block output
has the lower period required by
\eqref{fqi-period-threshold}. On each fixed starting torus
there are finitely many permitted steps. Uniformity as that
torus grows does not assert an infinite sequence on one
fixed finite regulator.

For representative selection, choose at stage $j$ a finite
family $E_j$ within $C_{\mathrm{int}}b^2$ of the quotient infimum.
Then
\begin{equation}
 \|E_j\|_{Q_j;a_0,b}
        \le2C_{\mathrm{int}}b^2<\bar\epsilon/2.
 \label{fqi-chosen-family}
\end{equation}
The infimum need not be attained. This family is holomorphic
on the full same $b$ ball, has connected carriers at the
same exponent, and has exactly the required normalized total.
Its total symmetries are unchanged. A real class permits the
norm-nonincreasing real averaging described earlier.

Use $E_j$ as the next input. Its norm is strictly inside the
actual family ceiling, so the inherited support estimates
and Gaussian theorem apply. Lemma~\ref{fqi-descent-lemma}
shows that this choice changes no next quotient output.
After the next step and actual finite reference subtraction,
choose another such family. Every generated correction remains
in the total. Since equality is exact, the positive
near-minimal tolerance creates no accumulating approximation
error and need not form a summable sequence. The raw constant
is handled by \eqref{fqi-vacuum-update} below.
\end{proof}

This selection occurs at each finite regulator; it is not
an inference from a small infinite class. Every winding
contribution is part of the finite total estimated by
Theorem~\ref{fcc-theorem} and the finite quantum bounds.
The assertion does not say that an untouched canonical
cluster array remains small. Selection is existential in
a finite Banach quotient, not a proposed algorithm for
finding a minimizing family. Exact equality and the uniform
norm are the properties needed for this induction.

The envelope also applies along an adaptively updated physical
coupling sequence for as long as it stays weak. For the bound
on a single such history one may freeze its realized coupling
values and use the auxiliary zero-started comparison with those
same values. This does not give a difference contraction
between histories that have different updated couplings.

\subsection{Exact vacuum compensation and coupling bookkeeping}
\label{fqi-vacuum-section}

Write the finite input density on the common chart as
\begin{equation}
 Z_j\exp[-\gamma_jS_j+E_j^{\rm raw}],\qquad
 \gamma_j=g_j^{-2}.
 \label{fqi-input-density}
\end{equation}
The finite exact chosen local step, with its known scalar
normalization $N_j$, produces
\begin{equation}
 Z_jN_j\exp[-\gamma_jS_{j+1}+F_j^{\rm raw}].
 \label{fqi-output-density}
\end{equation}
Let
\[
 \beta_{\rm out}=\beta_{j+1}([F_j^{\rm raw}]),\qquad
 x_{j+1}=[F_j^{\rm raw}]-\beta_{\rm out}[S_{j+1}],
\]
and choose the small family $E_{j+1}$ representing this class.
Since $S_{j+1}(0)=0$, equality of normalized totals gives
\begin{align}
 F_j^{\rm raw}(w)
   &=\beta_{\rm out}S_{j+1}(w)
                       +E_{j+1}^{\rm raw}(w)+c_j,\notag\\
 c_j&=F_j^{\rm raw}(0)-E_{j+1}^{\rm raw}(0).
 \label{fqi-constant-compensation}
\end{align}
The equality is on the entire common product domain. Hence
the exact updates preserving \eqref{fqi-output-density} are
\begin{equation}
 \gamma_{j+1}=\gamma_j-\beta_{\rm out},\qquad
 Z_{j+1}=Z_jN_j e^{c_j}.
 \label{fqi-vacuum-update}
\end{equation}
If an input representative is replaced by
$E'^{\rm raw}=E^{\rm raw}+c$ before integration, replace
$Z_j$ by $Z_je^{-c}$. Both signs are fixed by the raw
exponential, not by a convention about the quotient.
On real inputs use the compatible real logarithm branches;
constant $2\pi i$ changes in another compatible convention
do not change the exponentials. No convergence of an extensive
vacuum free energy is asserted.

If the couplings are freely prescribed without satisfying
\eqref{fqi-vacuum-update}, the theorem still bounds the
corresponding projected maps. It does not then assert
repeated factorization of one unchanged density. Along
the actual update, the amplitude bounds hold only while
the couplings stay in their proved weak range.

For the unprojected value along the envelope,
\begin{equation}
 \|\mathcal F_j(x_j)\|
 \le[B_TC_{\mathrm{int}}+C_F+L_{\mathrm{quant}}C_{\mathrm{int}}]b^2,
 \qquad L_{\mathrm{quant}}=\sup_j L_{\mathrm{quant},j}.
 \label{fqi-unprojected-bound}
\end{equation}
Separating $\mathcal T_jx_j$, $\mathcal G_j(0)$ and
$\mathcal G_j(x_j)-\mathcal G_j(0)$ proves this bound.
The finite Maxwell functional consequently gives
\begin{equation}
 |\gamma_{j+1}-\gamma_j|
 \le C_\beta[B_TC_{\mathrm{int}}+C_F
                              +L_{\mathrm{quant}}C_{\mathrm{int}}].
 \label{fqi-increment-bound}
\end{equation}
This is uniform in the finite regulator but has no sign.
For a specified finite number $n$ of steps, an initial
inverse coupling exceeding the weak threshold by $n$ times
this bound suffices for weak containment on that horizon.
It does not select a desired final coupling, establish an
averaged positive drift or control an infinite physical
trajectory. Comparing different running-coupling histories
would require an additional coupling-parameter comparison.

\subsection{Precise closure and remaining obligations}
\label{fqi-scope-section}

The theorem closes the finite \emph{representative-budget}
interface for this homogeneous local scheme. At every allowed
stage the chosen connected family has norm $O(b^2)$ at the
same carrier exponent and field radius. Every actual next
input and Gaussian budget holds; finite normalized-total
identities are exact; vacuum constants are compensated by
\eqref{fqi-constant-compensation}--\eqref{fqi-vacuum-update}.
The q-block amplitude contraction is at most one half for
equal prescribed coupling data.

This does not prove smallness of an untouched carrier array,
identify a one-shot composite sector with sequential primitive
sectors, supply large-field or mixed-history terms, or show
that the local germs represent every global Wilson history.
Those interpretations require their own orbit-domain and
measure identities. Every retained window remains precisely
the one in the common-domain theorem.

A physical density trajectory must obey
\eqref{fqi-vacuum-update} and remain weak. Its ultraviolet
tuning and beta drift, other field sectors, vacuum limits,
observables, nontriviality and a continuum measure remain
separate. None follows solely from the finite representative
induction proved here.


% No preamble. Prefix cgp-. Fixed T while complexifying g.
\section{Coupling parity and the leading zero-input correction}
\label{cgp-section}

Fix the composite geometry, auxiliary grouping, trace convention and
input-independent classical system of the homogeneous local macrostep.
Use one common outer field polydisc of radius $\rho>0$ and connected
carrier exponent $a_0$. All constants below may depend on that fixed
geometry and $\rho$, but are uniform in preceding homogeneous depth
and compatible volume. Infinite statements concern the same selected
local-limit system. No inherited interaction is present.

The zero-input correction has the exact decomposition
\begin{equation}
 \mathcal G_{M,g,T}(0;w)
   =\Delta_M(w)+\mathcal Q_{{\rm det},M}(w)
                         +\mathcal L_{M,g,T}(w),
 \qquad
 \mathcal L_{M,g,T}
  =\log\mathbb E_{C_M(w)}
          [\chi_T(\xi)e^{V_{\rm base}(w,\xi;g)}].
 \label{cgp-decomposition}
\end{equation}
The pivot term $\Delta_M$, Gaussian determinant term
$\mathcal Q_{{\rm det},M}$, covariance $C_M(w)$, stationary branch
and retained chart are independent of $g$. The actual cutoff
$\chi_T$ is a product of centered retained link/color balls.
Field-independent powers of $g$, flat determinants, Haar constants
and the compact gauge-insertion normalization belong to the
separately tracked vacuum factor. In particular no analyticity of
that exact compact normalization at $g=0$ is assumed.

We use the anchored norm
\[
 \|E\|_{a_0,\rho}
   =\sup_i\sum_{H\ni i}e^{a_0|H|}\|E_H\|_\rho
\]
and its quotient by equality of vacuum-normalized totals. A
coefficient family is not required to be individually gauge invariant.
The estimates below refer to this family norm or its indicated quotient,
not to an extensive scalar supremum.

\subsection{Exact parity of the centered factors}
\label{cgp-parity-section}

Keep $T$ fixed while allowing complex $g$. Every nonlinear
unwindowed factor is evaluated in the actual retained chart at the
same product $g\xi$. For a scalar action or gauge function $F$,
the centered contribution is
\begin{equation}
 -\frac1{g^2}\left[
 F(g\xi)-F(0)-g\,DF(0)\xi
                 -\frac{g^2}{2}D^2F(0)[\xi,\xi]\right].
 \label{cgp-centered-factor}
\end{equation}
The apparent singularity at zero is removable: the numerator
starts at total fluctuation degree three. Replacing $g$ by $-g$
is exactly the same as replacing $\xi$ by $-\xi$ in this
expression. This remains true on the independent centered
background extension; it does not rely on dropping a linear
term only at a physical stationary background.

The straightening determinant and Haar factors are analytic
functions of $g\xi$ with their near-origin logarithms, and obey
the same identity. They need not be separately even in $\xi$.
Consequently the exact total base exponent satisfies
\begin{equation}
             V_{\rm base}(w,\xi;-g)
                    =V_{\rm base}(w,-\xi;g).
 \label{cgp-base-parity}
\end{equation}
The centered Gaussian density is even in $\xi$, also for the
allowed complex background covariance. The fixed real window
is invariant under $\xi\mapsto-\xi$. This real change of
integration variable gives
\begin{equation}
                 Z_{M,-g,T}(w)=Z_{M,g,T}(w).
 \label{cgp-integral-parity}
\end{equation}
No contour deformation or change of the cutoff with the
background is involved.

\subsection{A complex disk uniform in the fixed threshold}
\label{cgp-disk-section}

The actual local chart has a common analytic neighborhood in
$g\xi$. The retained window bounds that variable by a fixed
constant times $|g|T$. The centered action, gauge and measure
estimates on the outer field domain give
\begin{equation}
 \|V_{\rm base}\|_{\rm strong}
   \le C_{\rm base}\bigl(|g|T^3+|g|T+|g|^2T^2\bigr).
 \label{cgp-base-budget}
\end{equation}
The strong norm is the one used by the strict Gaussian
support and overlap theorem; its constants are fixed before
$g,T$. Choose $T_0\ge1$ so the bad-block tail meets that
theorem's strict budget, and then choose a sufficiently
small fixed $c>0$. For all $T\ge T_0$, the disk
\begin{equation}
                         |g|<r_T:=cT^{-3}
 \label{cgp-disk}
\end{equation}
lies inside the local chart and the strict base budgets.
Indeed $|g|T\le cT^{-2}\le c$ and the right side of
\eqref{cgp-base-budget} is at most
$C_{\rm base}(2c+c^2)$. The precision and covariance
budgets do not move with $g$.

The strict Gaussian theorem therefore supplies a nonzero
integral and its holomorphic logarithm throughout
\eqref{cgp-disk}, with a connected coefficient family satisfying
\begin{equation}
             \sup_{|g|<r_T}
               \|\mathcal L_{M,g,T}\|_{a_0,\rho}
                         \le H_{\mathcal L}.
 \label{cgp-uniform-log}
\end{equation}
Here $H_{\mathcal L}$ is independent of $M$, volume and
$T\ge T_0$. Analyticity in $g$ uses the removable remainders
in \eqref{cgp-centered-factor}, analytic local measure factors,
finite Gaussian moments and the normally convergent positive
support/cluster majorants. The sharp indicator is fixed and
never differentiated.

The logarithm branch is fixed at $g=0$ by the same connected
Gaussian expansion. On the connected disk, the logarithms
in \eqref{cgp-integral-parity} differ by a constant multiple
of $2\pi i$, and that multiple is zero at $g=0$. Thus the
total logarithm is even.

\subsection{An even representative and the quantitative remainder}
\label{cgp-estimate-section}

\begin{theorem}[Fixed-threshold parity estimate]
\label{cgp-theorem}
Under the preceding actual-chart and strict-budget hypotheses,
for $T\ge T_0$ and $|g|T^3\le c/2$,
\begin{equation}
 \|[\mathcal L_{M,g,T}-\mathcal L_{M,0,T}]
                     \|_{\mathcal Q,a_0,\rho}
       \le 2H_{\mathcal L}
                       \left(\frac{|g|T^3}{c}\right)^2.
 \label{cgp-even-bound}
\end{equation}
Moreover, with
\begin{equation}
                D_M:=\Delta_M+\mathcal Q_{{\rm det},M},
 \label{cgp-leading-definition}
\end{equation}
there is a uniform constant $C_{\rm par}$ such that
\begin{equation}
 \|[\mathcal G_{M,g,T}(0)-D_M]\|_{\mathcal Q,a_0,\rho}
                 \le C_{\rm par}\bigl[\tau(T)+|g|^2T^6\bigr].
 \label{cgp-leading-bound}
\end{equation}
\end{theorem}
\begin{proof}
Evenness of a total does not require the original localized
coefficients to be even. Replace each coefficient by
\begin{equation}
 \mathcal L_H^{\rm ev}(g)
        :=\frac{\mathcal L_H(g)+\mathcal L_H(-g)}2.
 \label{cgp-even-representative}
\end{equation}
This retains its carrier and does not increase
\eqref{cgp-uniform-log}. Its total is the same logarithm.
Subtracting the coefficient family at $g=0$ represents the
actual difference in \eqref{cgp-even-bound}; that generally
background-dependent cutoff logarithm is not being declared
a vacuum constant.

Banach-valued Cauchy on radii increasing to $r_T$ gives
\begin{equation}
 \|\mathcal L^{\rm ev}(g)-\mathcal L^{\rm ev}(0)\|_{a_0,\rho}
 \le H_{\mathcal L}\sum_{n\ge1}(|g|/r_T)^{2n}
 \le2H_{\mathcal L}(|g|/r_T)^2.
 \label{cgp-Cauchy-proof}
\end{equation}
The same estimate follows by taking the Cauchy integral before
the positive anchored sums. It uses no interchange of an
order sum with a root supremum and no count of background
coordinates. Taking the quotient proves
\eqref{cgp-even-bound}.

At $g=0$ all centered action and gauge remainders and the
$g\xi$ measure remainders vanish. The remaining logarithm is
precisely the Gaussian cutoff logarithm, whose strict
connected norm is at most $C\tau(T)$. Combine this with
\eqref{cgp-even-bound} and \eqref{cgp-decomposition}.
\end{proof}

\subsection{The physical threshold and the Maxwell functional}
\label{cgp-Maxwell-section}

Only after the fixed-$T$ estimate has been proved do we evaluate
it at the prescribed real threshold $T=T(g)$. This is an
inequality in two parameters, not a claim that the sharp
$g$-dependent window is holomorphic. For a fixed sufficiently
large multiple of $\sqrt{\log(1/g)}$ and a fixed lower
threshold,
\[
 \tau(T(g))\longrightarrow0,\qquad
          g^2T(g)^6=O\!\left(g^2\log(1/g)^3\right).
\]
The tail can be made a prescribed positive power of $g$
by its fixed coefficient, subject to the stated Gaussian
tail bound.

The fixed-order kernel estimate on the outer field domain gives
$\|\beta\|\le C_\beta\rho^{-2}$. Therefore
\begin{equation}
 \left|\beta(\mathcal G_{M,g,T(g)}(0))-\beta(D_M)\right|
 \le C_\beta\rho^{-2}C_{\rm par}
                \bigl[\tau(T(g))+g^2T(g)^6\bigr].
 \label{cgp-beta-bound}
\end{equation}
This convergence is uniform in $M$ under the established
common bounds. It neither computes $\beta(D_M)$ nor proves
that those coefficients converge as $M$ grows.

There is also a useful inner-radius form for other bounded
linear marginal coordinates. The actual zero-input totals
are locally gauge and globally adjoint invariant.
Equation \eqref{cgp-leading-bound} along $T(g)$ shows that
$D_M$ is an invariant total as well, since the invariant
subspace of the outer quotient is closed. For a semisimple
Lie algebra the total first derivative of their difference
vanishes. Remove the degree-zero and degree-one background
sectors of a representative of that difference and restrict
to $b\le\rho/2$. Homogeneous phase projection gives
\begin{equation}
 \|[\mathcal G_{M,g,T(g)}(0)-D_M]\|_{\mathcal Q,a_0,b}
 \le \frac{2C_{\rm par}}{\rho^2}\,
       b^2\bigl[\tau(T(g))+g^2T(g)^6\bigr].
 \label{cgp-inner-bound}
\end{equation}
Taking representatives arbitrarily close to the outer
quotient infimum justifies the same assertion when the
outer estimate was stated only as a quotient bound.
No coefficientwise gauge identity is needed.

\subsection{Scope}
\label{cgp-scope-section}

The terms $\Delta_M$ and $\mathcal Q_{{\rm det},M}$ have no
small factor $g$. Their removal of total constant and linear
sectors gives a quadratic inner-field-radius bound, not an
$O(g^2)$ bound for the full correction.
The theorem identifies a concrete leading determinant/pivot
functional and a vanishing error. It proves no sign for its
Maxwell coefficient, no running-coupling trajectory and no
identification of the local sector with the full Wilson measure.
All vacuum factors remain separately tracked. The fixed
classical-system and common-domain qualifications are unchanged.


% No preamble. Prefix cmt-. True infinite Maxwell coefficient; not finite cutoff beta.
\section{Exact classical preservation of the Maxwell coefficient}
\label{cmt-section}

This section concerns the invariant infinite interaction quotient
on one selected joint classical background system. The coefficient
\(\beta_M\) is the true infinite total-kernel Maxwell moment.
It is not the finite symmetric-cutoff ratio \(\beta_{M,N}\).
Fix the primitive ratio, grouping \(J\), activity exponent \(a_0\),
and one positive field radius \(b\). The actual homogeneous
reference and projection satisfy
\begin{equation}
 \beta_M(S_M)=1,\qquad \|S_M\|_b\le C_Sb^2,\qquad
 P_M=I-S_M\beta_M,\qquad \|P_M\|\le C_{\rm ext}(J).
 \label{cmt-reference}
\end{equation}
The actual pre-gauge substitution has the scalar identities
\begin{equation}
 \mathcal T_{M,L}S_M=S_{ML},\qquad
 \mathcal T_{MR,L}\mathcal T_{M,R}
                       =\mathcal T_{M,RL}.
 \label{cmt-semigroup}
\end{equation}
They are identities of normalized interactions, not of chosen
support coefficient arrays.

The complete-field pre-gauge family obeys, for every fixed
prescribed finite classical reserve \(A\),
\begin{equation}
 \sup_i\sum_X e^{\kappa\tau(\{i\}\cup X)+A|X|}
       \|a_{M,L;i,X}\|_{C^{1,\alpha};b}
           \le C_* b/L,\qquad 0<b\le r_A,
 \label{cmt-uniform-family}
\end{equation}
on one outer radius \(r_A>0\) independent of \(M,L\).
The entire output fiber and its physical first-gradient
H\"older quotients are inside the support sum, and the
coefficients agree on overlapping collars. Fix \(J\) and
the reserves so that \(2(a_0+1)/J\le\kappa\) and
\(A\ge4(a_0+1)\).

The argument uses countably many arbitrarily large endpoint
powers. They can be included in the chosen joint system by
a further diagonal extraction with the same \(r_A\);
previously fixed limits remain unchanged. The exact finite
comparison and controlled gauge identities then give
\eqref{cmt-semigroup} on this enlarged joint system.
No independence from unrelated subsequences is asserted.

\subsection{The quadratic estimate needs no large field-scaling disk}
\label{cmt-quadratic-section}

Let \(y=P_Mx\) for an invariant class \(x\). Global semisimple
color invariance gives \(D\Phi_y(0)=0\). The exact infinite Ward
classification and hypercubic symmetry give
\begin{equation}
 \Gamma_2^y(p)=O(|p|^4),\qquad
 \|K_2(y)\|_{\sigma(J),{\rm row}}
                         \le C(J)b^{-2}\|y\|_b.
 \label{cmt-kernel}
\end{equation}
The fixed midpoint phases preserve the order of vanishing in
tail coordinates. Define \(Q_y(A)=K_2(y)[A,A]/2\).
For compact coarse \(w\), the pre-gauge image is exponentially
decaying and summable, so \(Q_y(a_{M,L}(w))\) is well defined.

The quadratic part of the spatial localization proof yields
a connected family \(Q_{y,L}\) representing that entire
substituted quadratic polynomial, with
\begin{equation}
 \|Q_{y,L}\|_{a_0,b}
       \le C_2(J)L^{-\alpha}\|y\|_b,
       \qquad L\ge L_0(J),\quad 0<b\le r_A.
 \label{cmt-quadratic-bound}
\end{equation}
Indeed the divided-difference factorization is
\(K_2=d^-\widetilde H d^+\), with
\(\sum_r\widetilde H(r)=0\), and gives the exact expression
\[
 \frac12\sum_{x,r,\rho,\tau}
 [D_\rho a(x)]^T\widetilde H_{\rho\tau}(r)
                    [D_\tau a(x+r)-D_\tau a(x)].
\]
The complete family bounds supply \(CbL^{-2}\) for the first
difference and \(CbL^{-2-\alpha}|r|^\alpha\) for the nearby
gradient difference. The exponential kernel tail controls
the far terms. The \((JL)^4\) anchor count leaves
\(L^{-\alpha}\). All constants are fixed after \(J\) and
before \(L\).

Only \(a_{M,L}(w)\) on the original \(b\)-ball is used.
There is no substitution \(a_{M,L}(\lambda w)\) with
\(|\lambda|\) proportional to \(L\). Thus the high-jet
condition \(bL\le\text{constant}\) is absent from
\eqref{cmt-quadratic-bound}.

Although \(Q_{y,L}\) need not be invariant as a whole nonlinear
functional, the chosen Hessian moment extends to every analytic
interaction quotient as a bounded linear functional:
\begin{equation}
 |b_{\rm mom}(Z)|\le C_\beta(J)b^{-2}\|Z\|_b.
 \label{cmt-moment-bound}
\end{equation}
On invariant classes it is \(\beta\). This norm bound is just
a fixed second derivative and a fixed second spatial moment;
it does not require a Ward identity for \(Z\).

\subsection{The transported Hessian and the semigroup limit}
\label{cmt-preservation-section}

Since \(a_{M,L}(0)=0\) and \(D\Phi_y(0)=0\), the chain rule gives
\begin{equation}
 D^2(\Phi_y\circ a_{M,L})(0)[u,v]
   =D^2\Phi_y(0)[Da_{M,L}(0)u,Da_{M,L}(0)v].
 \label{cmt-chain-rule}
\end{equation}
The right side is also the Hessian of
\(Q_y\circ a_{M,L}\). Higher input field degrees cannot
contribute to degree two, and the possible
\(D\Phi_y(0)D^2a_{M,L}(0)\) term is zero.
Differentiation is justified first on finitely supported
directions by local normal convergence, then for the moment
by the anchored exponential kernel bounds. Hence
\begin{equation}
 \beta_{ML}(\mathcal T_{M,L}y)=b_{\rm mom}(Q_{y,L}).
 \label{cmt-hessian-identification}
\end{equation}
Combining \eqref{cmt-quadratic-bound} and
\eqref{cmt-moment-bound}, followed by the exact reference
transport, proves
\begin{equation}
 |\beta_{ML}(\mathcal T_{M,L}x)-\beta_M(x)|
       \le C_m(J)b^{-2}L^{-\alpha}\|x\|_b.
 \label{cmt-asymptotic}
\end{equation}
This estimate is uniform in preceding depth, at fixed \(b>0\).

\begin{theorem}[Exact infinite Maxwell preservation]
\label{cmt-theorem}
For every fixed admissible ratio \(R\) and every invariant
class \(x\) on this common domain,
\begin{equation}
       \beta_{MR}(\mathcal T_{M,R}x)=\beta_M(x).
 \label{cmt-exact}
\end{equation}
\end{theorem}
\begin{proof}
Put \(z=\mathcal T_{M,R}x\). The fixed-\(R\) classical
transport is bounded, say
\(\|z\|_b\le B_R(J)\|x\|_b\). Apply
\eqref{cmt-asymptotic} to \(z\) at depth \(MR\) and
endpoint \(L\), and to \(x\) at depth \(M\) and endpoint
\(RL\). By \eqref{cmt-semigroup}, their leading transported
moments agree. Therefore
\begin{equation}
 |\beta_{MR}(\mathcal T_{M,R}x)-\beta_M(x)|
 \le C_m(J)b^{-2}
       \bigl[B_R(J)L^{-\alpha}+(RL)^{-\alpha}\bigr]\|x\|_b.
 \label{cmt-comparison}
\end{equation}
The left side is independent of \(L\). Let \(L\) tend through
the admissible endpoint powers to infinity in the joint
system.
\end{proof}

The proof needs no convergence of the full interactions
\(\mathcal T_{M,L}x\) as \(L\to\infty\), and no bound on
their whole norms uniform in \(L\). Only the fixed-\(R\)
map, the scalar semigroup, and the quadratic-only bound
are used.

There is no hidden domain shrink. The actual pre-gauge
holomorphy and \eqref{cmt-uniform-family} have a fixed
primitive-scale radius. A real minimum comparison at
one large endpoint can initially require a smaller real
germ because physical norms rescale. Both scalar sides
of \eqref{cmt-semigroup}, however, are holomorphic on
their specified common \(b\)-domain; equality on that
nonempty germ extends throughout the connected domain.
On compact field subspaces the anchored and summable-field
bounds justify that continuation. The high-jet dilation
radius plays no role. Countably many endpoints can
therefore be included without replacing \(b\) by a
sequence tending to zero.

\subsection{The original Maxwell coordinate isolates the leading drift}
\label{cmt-drift-section}

Equations \eqref{cmt-exact} and \eqref{cmt-semigroup} imply
\begin{equation}
 \beta_{MR}\mathcal T_{M,R}=\beta_M,\qquad
 \mathcal T_{M,R}P_M=P_{MR}\mathcal T_{M,R}.
 \label{cmt-commutation}
\end{equation}
Thus the classical leakage on \(\ker\beta_M\) is exactly
zero for this infinite system. An adapted-coordinate
correction introduced solely to remove that leakage
can be chosen zero; the general coordinate theorem
remains valid for other operator families.

For the fixed-coupling quantum map
\(\mathcal F_{M,g,T}=\mathcal T_{M,R}+\mathcal G_{M,g,T}\),
let \(x=P_Mx\), \(\|x\|_b\le C_{\rm int}b^2\). The parity
theorem defines
\(D_M=\Delta_M+\mathcal Q_{{\rm det},M}\), as in
\eqref{cgp-leading-definition}, and gives the inner bound
\eqref{cgp-inner-bound}. The input Lipschitz estimate gives
\[
 \|\mathcal G_{M,g,T}(x)-\mathcal G_{M,g,T}(0)\|_b
      \le C_{\rm quant}(gT/b)C_{\rm int}b^2.
\]
Applying the original Maxwell moment, with
\(\|\beta_{MR}\|\le C_\beta(J)b^{-2}\), yields
\begin{align}
 &|\beta_{MR}(\mathcal F_{M,g,T}(x))
                         -\beta_{MR}(D_M)|\notag\\
 &\quad\le
   \frac{2C_\beta C_{\rm par}}{\rho^2}
                         [\tau(T)+g^2T^6]
       +C_\beta C_{\rm quant}C_{\rm int}\,gT/b.
 \label{cmt-drift-bound}
\end{align}
All constants are uniform in preceding depth on the fixed
classical system. The field radius \(b\) is fixed before
the real threshold \(T=T(g)\) and \(g\to0\) are considered.
For an actual density update
\(\gamma_{\rm next}=\gamma-\beta_{MR}(\mathcal F(x))\),
the leading decrement in the original coupling coordinate
is therefore \(\beta_{MR}(D_M)\), with the displayed
vanishing error. No adapted marginal coordinate is needed
for this conclusion.

\subsection{Finite and physical limitations}
\label{cmt-scope}

The finite coefficient \(\beta_{M,N}\) is a short-lift
ratio and is not the coefficient in
\eqref{cmt-exact}. Its Ward defects and folded tails
are nonzero at finite period. This argument does not
prove exact finite marginal preservation or set the
finite classical leakage to zero. The finite reference
telescope and finite contraction remain valid without
that identity.

Equation \eqref{cmt-drift-bound} does not compute the
determinant/pivot coefficient, prove its sign or convergence
with \(M\), or control an infinite physical coupling
trajectory. The result concerns the selected joint
infinite classical system and does not establish
subsequence independence, global sector matching,
observable convergence or continuum existence.


\section{The coupling calculation and the remaining existence argument}
\label{fin-next}

The finite local induction supplies actual admissible interaction
families at every stage. Its exact density bookkeeping is
\[
 \gamma_{j+1}=\gamma_j-
      \beta_{j+1}\bigl(\mathcal F_{j,g_j}(E_j)\bigr),
 \qquad \gamma_j=g_j^{-2}.
\]
A trajectory of one physical density must obey this equation.
The bounded magnitude of its increments does not determine their
sign or tune an arbitrarily long ultraviolet sequence to a fixed
interacting endpoint.

The parity theorem isolates a concrete leading functional,
\[
 D_M=\Delta_M+\mathcal Q_{{\rm det},M}.
\]
Exact infinite classical transport and the quantum input bound
reduce the drift on the \(O(b^2)\) amplitude envelope to
\[
 \beta_{j+1}(D_{M_j})+
 O\!\left(\tau(T(g_j))+g_j^2T(g_j)^6+
                         \frac{g_jT(g_j)}b\right).
\]
Here \(b>0\) is fixed before sending the coupling to zero;
constants may depend on that fixed field radius and geometry.
This expression concerns the infinite quotient, with the
finite cutoff-extraction distinction retained above.
The leading Maxwell coefficient must be evaluated using the actual
stationary background, pivot coordinates, gauge normalization and
retained Gaussian. Every implicit background derivative belongs
to that calculation. A universal perturbative coefficient from
another scheme cannot substitute for this identification.

The next measure-theoretic steps remain substantial. The homogeneous
local germ must be represented on the required gauge-orbit domains,
and mixed and large-field histories must be incorporated with their
exact measure identities. Uniform source estimates must then produce
a nontrivial continuum observable hierarchy with the symmetry,
positivity and reconstruction properties specified in Notebook 01.
No conclusion about those steps follows merely from changing
interaction representatives or bounding a finite local integral.

\begin{thebibliography}{9}
\bibitem{priorfin}
\emph{Toward a four-dimensional Yang--Mills construction},
research Notebooks 15--30 in this workspace. The finite ordinary
minimum, common-domain local quantum map, total Ward extraction
and quantum quotient estimates are the principal inputs.
\bibitem{balaban109fin}
T.~Ba\l aban,
\emph{Renormalization group approach to lattice gauge field
theories. I. Generation of effective actions in a small field
approximation and a coupling constant renormalization in four
dimensions},
Commun. Math. Phys. \textbf{109} (1987), 249--301.
\href{https://doi.org/10.1007/BF01215223}{Publisher record}.
\bibitem{balaban116fin}
T.~Ba\l aban,
\emph{Renormalization group approach to lattice gauge field
theories. II. Cluster expansions},
Commun. Math. Phys. \textbf{116} (1988), 1--22.
\href{https://doi.org/10.1007/BF01239022}{Publisher record}.
The finite quotient and representative arguments here are given
separately; the published induction theorem is not being imported
with altered hypotheses.
\bibitem{sourcefin}
E.~Sussman and collaborators,
\href{https://github.com/ethanwsussman/Yang--Mills-LLM-experiment}
{Yang--Mills-LLM-experiment}, deposited raw accounts at commit
\texttt{e392ef45013a581664bab76b90867170b30fe10e}.
\end{thebibliography}
\end{document}
